\documentclass{article}
\usepackage{/globals/de}
\usepackage{/globals/math}
\usepackage{/globals/boxes}
\usepackage{/globals/anki}
\ankiprefix{bases-decimal}
\ankitopic{Zahlensysteme Kommazahlen}
\ankishow
\begin{document}
\section{Das Hornerschema} 
Mit dem Hornerschema können Zahlen aus der Basis $B$ in das Dezimalsystem konvertiert werden ohne dass mehrere Potenzen berechnet werden müssen. Stattdessen wird während der einen Berechnung selbst mehrfach multipliziert oder Dividiert.

\subsection{Ganze Zahlen und Vorkommastellen}
Für den Vorkommateil, und für ganze Zahlen, gilt
\[
 n_{10} = ((b_{k-1}) * B + b_{k-2}) * B \dots + b_0
\]
Wegen Punkt vor Strich Regeln und den vielen Klammern sieht diese Rechnung unnötig komplex aus; es wird einfach wiederholt von links nach rechts eine Ziffer addiert und alles mit der Basis multipliziert, die nächste Ziffer wird addiert, wieder alles mit der Basis multipliziert, bis die letzte Ziffer addiert wird.
\begin{zexample}
Es wird $1101_2$ in das Dezimalsystem umgerechnet.
\begin{align*}
   &((1*2+1)*2+0)*2+1 \\
=\ &3*2*2+1 \\
=\ &13
\end{align*}
Oder auch, wird jeder Schritt nacheinander aufgeschrieben
\begin{align*}
 0\ &\vert\ (+1) \\
 1\ &\vert\ (*2) \\
 2\ &\vert\ (+1) \\
 3\ &\vert\ (*2) \\
 6\ &\vert\ (+0) \\
 6\ &\vert\ (*2) \\
 12\ &\vert\ (+1) \\
 13\ &
\end{align*}
Somit ist $1101_2 = 13_{10}$.
\end{zexample}

\subsection{Nachkommastellen}
Für den Nachkommateil gilt dem ähnlich
\[
 n_{10} = ((b_{m-1}) * \frac{1}{B} + b_{m-2}) * \frac{1}{B}) \dots + b_{-1}
\]
Die Idee ist die Gleiche, nur dass hier die Ziffern von rechts nach links gelesen werden, mit $1/B$ multipliziert, \bzw durch $B$ dividiert wird, und dass es eine weitere multiplikation am Ende gibt.
\begin{zexample}
Es wird $0,1101_2$ in das Dezimalsystem umgerechnet.
\begin{align*}
   &\left(\left(\left(1*\frac{1}{2}+0\right)*\frac{1}{2}+1\right)*\frac{1}{2}+1\right)*\frac{1}{2} \\
=\ &0,8125
\end{align*}
Oder auch, wird jeder Schritt nacheinander aufgeschrieben
\begin{align*}
 0\ &\vert\ (+1) \\
 1\ &\vert\ (/2) \\
 0,5\ &\vert\ (+0) \\
 0,5\ &\vert\ (/2) \\
 0,25\ &\vert\ (+1) \\
 1,25\ &\vert\ (/2) \\
 0,625\ &\vert\ (+1) \\
 1,625\ &\vert\ (/2) \\
 0,8125\ &
\end{align*}
Somit ist $0,1101_2 = 0,8125_{10}$.
\end{zexample}
\end{document}